\chapter{Multi-View Geometry}
\label{ch:multiview}

\begin{objectives}
\begin{itemize}
  \item Derive the epipolar constraint and the essential matrix.
  \item Estimate relative motion from correspondences with the 8-, 5- and 2-point methods.
  \item Use RANSAC to reject outliers, size its iteration count, and set thresholds in the right units.
  \item Understand the pure-rotation degeneracy and how to handle it.
  \item Triangulate points linearly and by Gauss--Newton with an inverse-depth parameterization, and quantify
        triangulation accuracy as a function of parallax.
\end{itemize}
\end{objectives}

% ------------------------------------------------------------------------------
\section{Epipolar geometry}

Two cameras observe the same 3D point. Let the coordinates of the point in the two camera frames be $X_1$ and
$X_2$, related by the rigid motion
\begin{equation}
  X_2 = R\,X_1 + t, \qquad R = R_{c_2c_1},\quad t = p_{c_2c_1},
\end{equation}
and let $\bar x_1 = X_1/Z_1$ and $\bar x_2 = X_2/Z_2$ be the bearings (homogeneous normalized coordinates).

\begin{theorem}[Epipolar constraint]
\label{thm:epipolar}
Every correspondence of a rigid point satisfies
\begin{equation}
  \bar x_2\T\,E\,\bar x_1 = 0, \qquad E = \sx{t}\,R .
  \label{eq:epipolar}
\end{equation}
\end{theorem}
\begin{proof}
Cross both sides of $X_2 = RX_1 + t$ with $t$: $\sx{t}X_2 = \sx{t}RX_1$. The left side is orthogonal to $X_2$, so
$X_2\T\sx{t}RX_1 = 0$. Dividing by $Z_1Z_2 > 0$ gives \eqref{eq:epipolar}.
\end{proof}

Geometrically, the two camera centers and the 3D point span a plane, the \emph{epipolar plane}, and the
constraint states that the two viewing rays and the baseline are coplanar (\cref{fig:epipolar}). Given $\bar x_1$, the
corresponding point in the second image must lie on the \emph{epipolar line} $\ell_2 = E\bar x_1$. The constraint involves no
depth: it is a statement about bearings only, which is why it is the natural tool for checking tracks before any 3D
reconstruction exists.

\begin{figure}[t]
\centering
\begin{tikzpicture}[scale=1.0]
  \coordinate (C1) at (0,0); \coordinate (C2) at (5,0.3); \coordinate (P) at (2.7,3.4);
  \fill[primarylight, opacity=0.7] (C1) -- (C2) -- (P) -- cycle;
  \draw[thick] (C1) -- (C2) node[midway, below] {baseline $t$};
  \draw[plotgreen, thick] (C1) -- (P); \draw[plotgreen, thick] (C2) -- (P);
  \fill (C1) circle (1.8pt) node[below left] {$C_1$};
  \fill (C2) circle (1.8pt) node[below right] {$C_2$};
  \fill[conv] (P) circle (2pt) node[above] {$p_f$};
  % image planes
  \draw[primary, thick] (0.3,1.5) -- (1.6,1.2) -- (1.6,0.2) -- (0.3,0.5) -- cycle;
  \draw[primary, thick] (3.5,1.3) -- (4.8,1.7) -- (4.8,0.7) -- (3.5,0.3) -- cycle;
  \fill[plotgreen] ($(C1)!0.3!(P)$) circle (1.5pt) node[left, font=\footnotesize] {$\bar x_1$};
  \fill[plotgreen] ($(C2)!0.3!(P)$) circle (1.5pt) node[right, font=\footnotesize] {$\bar x_2$};
  \draw[conv, dashed] (3.6,0.65) -- (4.75,1.45);
  \node[conv, font=\footnotesize, anchor=west] at (5.0,1.25) {epipolar line $E\bar x_1$};
  \node[font=\footnotesize, primary] at (2.6,1.7) {epipolar plane};
\end{tikzpicture}
\caption{Epipolar geometry. The camera centers, the point and both viewing rays lie in one plane. The image of the
point in the second camera must therefore lie on the line where that plane cuts the second image.}
\label{fig:epipolar}
\end{figure}

\begin{proposition}[Properties of the essential matrix]
$E = \sx{t}R$ has rank two (for $t \neq 0$), two equal non-zero singular values, and five degrees of freedom: three
for the rotation and two for the direction of $t$. Its overall scale, like the scale of $t$, is unobservable.
\end{proposition}

In pixel coordinates the same constraint reads $\tilde u_2\T F\tilde u_1 = 0$ with the \emph{fundamental matrix}
$F = K_2\mT EK_1\inv$. The project works with undistorted normalized coordinates throughout, so it uses $E$ directly
(Milestone 9).

% ------------------------------------------------------------------------------
\section{Estimating relative motion}

\paragraph{The eight-point algorithm.} \Cref{eq:epipolar} is linear in the nine entries of $E$. Each correspondence gives
one equation $a_i\T\vect(E) = 0$, where $a_i$ contains the products of the coordinates of $\bar x_1$ and $\bar x_2$. With eight
or more correspondences, stack them into $A\,\vect(E) = 0$ and take the right singular vector of $A$ with the smallest
singular value. The result is then projected onto the set of valid essential matrices by setting its singular values
to $(1, 1, 0)$. Finally $R$ and the direction of $t$ are extracted by a second SVD; of the four candidate solutions, the
one that places triangulated points in front of both cameras is kept.

\paragraph{The five-point algorithm.} Five correspondences are the minimum, since $E$ has five degrees of freedom.
Nistér's algorithm~\cite{nister2004five} solves the resulting polynomial system and returns up to ten candidates. Fewer
points per hypothesis means fewer RANSAC iterations, which is why \code{cv::findEssentialMat} uses it.

\paragraph{The two-point algorithm with known rotation.} In VIO the gyroscope supplies the relative rotation $R$ between
consecutive frames with high accuracy. With $R$ known, \eqref{eq:epipolar} becomes linear in $t$:
\begin{equation}
  \bar x_2\T\sx{t}R\bar x_1 = t\T\big(R\bar x_1\times\bar x_2\big) = 0 .
\end{equation}
Each correspondence says that $t$ is orthogonal to the vector $c_i = R\bar x_1^{(i)}\times\bar x_2^{(i)}$, so two correspondences
determine the direction of $t$ as $c_1\times c_2$. A two-point hypothesis needs far fewer RANSAC iterations than a five-point
one, and it is the outlier-rejection strategy of MSCKF-VIO~\cite{sun2018msckfvio}.

\paragraph{Inlier distance.} The algebraic residual $\bar x_2\T E\bar x_1$ has no geometric meaning. The \emph{Sampson distance}
\begin{equation}
  d_S = \frac{(\bar x_2\T E\bar x_1)^2}{(E\bar x_1)_1^2 + (E\bar x_1)_2^2 + (E\T\bar x_2)_1^2 + (E\T\bar x_2)_2^2}
\end{equation}
is a first-order approximation of the squared reprojection distance, in the units of the coordinates used. With
normalized coordinates, a pixel threshold $\tau_{px}$ must be converted to $\tau_{px}/f$ (Milestone 9).

\begin{pitfall}[thresholds in the wrong units]
Calling \code{cv::findEssentialMat} on normalized coordinates with an identity camera matrix and a threshold of
``1.0'' accepts residuals of 1 normalized unit, about 450 pixels: every track becomes an inlier. The threshold
must be $\tau_{px}/f$, e.g.\ $1/458 \approx 2.2\times10^{-3}$.
\end{pitfall}

% ------------------------------------------------------------------------------
\section{RANSAC}

Random Sample Consensus~\cite{fischler1981ransac} estimates a model from data containing outliers by hypothesis and
verification.

\begin{algorithm}[h]
\caption{RANSAC for relative motion}
\begin{algorithmic}[1]
\Require correspondences, minimal sample size $s$, threshold $\tau$, confidence $p$
\State best $\gets \emptyset$
\For{$k = 1, \dots, N$}
  \State draw $s$ correspondences at random; fit candidate models
  \State count inliers with distance $< \tau$; keep the model with the most inliers
  \State update $N$ from the current inlier ratio $w$ with \eqref{eq:ransac-n}
\EndFor
\State refit the model on all inliers of the best hypothesis
\end{algorithmic}
\end{algorithm}

If a fraction $w$ of the data are inliers, a random minimal sample is outlier-free with probability $w^s$. To draw at
least one clean sample with probability $p$, the number of iterations must be
\begin{equation}
  N = \frac{\log(1 - p)}{\log(1 - w^s)} .
  \label{eq:ransac-n}
\end{equation}
\begin{center}\small
\begin{tabular}{@{}l c c c@{}}
\toprule
$p = 0.99$ & $w = 0.9$ & $w = 0.7$ & $w = 0.5$ \\
\midrule
8-point ($s = 8$) & 9 & 78 & 1177 \\
5-point ($s = 5$) & 6 & 26 & 146 \\
2-point ($s = 2$) & 3 & 7 & 17 \\
\bottomrule
\end{tabular}
\end{center}
The exponential dependence on $s$ is why minimal solvers matter, and why using the gyroscope to remove three
unknowns pays off.

% ------------------------------------------------------------------------------
\section{Degeneracies}

\subsection{Pure rotation}

If the camera only rotates, $t = 0$ and $E = \sx{0}R = 0$: the epipolar constraint is satisfied by every pair of points,
inlier or not. Under small translation the problem is not degenerate but badly conditioned: $t$ is poorly determined,
inlier sets fluctuate, and RANSAC may reject good tracks or accept bad ones. Hovering drones rotate in place often,
so this case must be handled (Milestone 9).

Under pure rotation the correspondences satisfy a stronger constraint, $\bar x_2 \propto R\bar x_1$ (the \emph{infinite homography}),
which does not involve $t$ at all. A robust front end measures the parallax remaining after rotation compensation,
for example the median of $\|\pi(R\bar x_1) - \pi(\bar x_2)\|$ over the tracks:
\begin{itemize}
  \item if it is large, run essential-matrix (or two-point) RANSAC;
  \item if it is tiny, skip epipolar RANSAC and reject only tracks that violate the rotation-only prediction
        $\bar x_2 \propto R\bar x_1$, with $R$ from the gyroscope.
\end{itemize}
Either way the system keeps tracking: low parallax means the features carry little information about translation,
not that they are wrong. The MSCKF update handles them gracefully, because their nullspace-projected residuals
simply contribute little.

\subsection{Other degeneracies}

The eight-point algorithm is degenerate when all points lie on a plane, a common situation (floors, walls); the
five-point algorithm is not. Very distant points behave like points at infinity and carry only rotational
information. Moving objects form a second, consistent rigid motion that RANSAC can mistake for the dominant one if they
fill the image.

% ------------------------------------------------------------------------------
\section{Triangulation}
\label{sec:triangulation}

Given observations $z_i$ of a point in cameras with known poses $(R_{wc_i}, p_{wc_i})$, triangulation estimates its world
position $p_f$.

\subsection{Linear triangulation}

In camera $i$ the point has coordinates $X_i = R_{wc_i}\T(p_f - p_{wc_i})$, which must be parallel to the bearing
$\bar x_i = (z_i, 1)$: $\sx{\bar x_i}X_i = 0$. Substituting,
\begin{equation}
  \sx{\bar x_i}R_{wc_i}\T\, p_f = \sx{\bar x_i}R_{wc_i}\T\, p_{wc_i}, \qquad i = 1,\dots,M .
  \label{eq:dlt}
\end{equation}
Each view contributes three equations of rank two. Stacking all views gives an overdetermined linear system
$A p_f = b$ that is solved in the least-squares sense with QR or SVD. With two views and noise-free data the solution is
exact. This is the linear initialization of AGENTS.md F5. It minimizes an algebraic error, not the reprojection error,
so it is refined next.

\subsection{Nonlinear refinement with inverse depth}

The maximum-likelihood estimate minimizes the reprojection error
\begin{equation}
  \min_{p_f}\ \sum_{i=1}^M \big\| z_i - \pi(X_i(p_f)) \big\|^2 .
\end{equation}
The MSCKF~\cite{mourikis2007msckf} parameterizes the point in an \emph{anchor} camera $A$ (for example the first observation)
by inverse depth~\cite{civera2008inverse}:
\begin{equation}
  X_A = \frac1\rho\begin{bmatrix} \alpha \\ \beta \\ 1 \end{bmatrix}, \qquad
  X_i = R_{iA}X_A + p_{iA}
  \quad\Longrightarrow\quad
  h_i := \rho X_i = R_{iA}\begin{bmatrix}\alpha\\ \beta\\ 1\end{bmatrix} + \rho\,p_{iA}, \quad z_i = \pi(h_i),
\end{equation}
where $(R_{iA}, p_{iA})$ is the transformation from the anchor camera to camera $i$, computed from the clone poses, and
$\pi(h_i) = \pi(X_i)$ because projection ignores scale. The parameters $(\alpha, \beta)$ are the normalized coordinates in
the anchor view and $\rho = 1/Z_A$ is the inverse depth. This parameterization behaves well for distant points
($\rho\to0$ is a regular point), and the measurement function is nearly linear in $(\alpha, \beta, \rho)$. Its Jacobian
follows from the chain rule:
\begin{equation}
  \frac{\partial z_i}{\partial(\alpha, \beta, \rho)} = \frac{\partial\pi}{\partial h}\Big|_{h_i}\,
  \begin{bmatrix} r_1 & r_2 & p_{iA} \end{bmatrix},
\end{equation}
where $r_1, r_2$ are the first two columns of $R_{iA}$ and $\partial\pi/\partial h$ is \cref{prop:proj-jacobian}. Gauss--Newton
iterates $\delta y = -(J\T J)\inv J\T e$ with the stacked Jacobian $J$ and residual $e$, usually converging in a few
iterations from the linear initialization. Levenberg--Marquardt damping adds robustness.

\subsection{How accurate is triangulation?}

For two views separated by a baseline $b$ perpendicular to the viewing direction, the depth is $Z = b/d$ with $d$ the
disparity in normalized coordinates. With independent noise $\sigma_{uv}$ in each view, $\sigma_d = \sqrt2\sigma_{uv}$ and to first
order
\begin{equation}
  \sigma_Z \approx \frac{Z^2}{b}\sqrt2\,\sigma_{uv} .
  \label{eq:depth-sigma}
\end{equation}
Depth error grows quadratically with depth and inversely with baseline. Equivalently, the relative depth error is
inversely proportional to the parallax angle $\approx b/Z$.

\begin{figure}[t]
  \centering
  \includegraphics[width=0.86\textwidth]{fig_triangulation.pdf}
  \caption{Relative depth error of two-view triangulation at $Z = 5$~m with 1-pixel noise, versus parallax angle.
  Above about $2^\circ$ the first-order formula \eqref{eq:depth-sigma} is accurate. Below about $1^\circ$ the error distribution
  becomes heavy-tailed (depths can even become negative), and the estimate is unusable.}
  \label{fig:triangulation}
\end{figure}

\Cref{fig:triangulation} shows that 10\% depth accuracy needs about $2^\circ$ of parallax, a baseline of 17~cm at 5~m. Below
about $1^\circ$ the linearization behind \eqref{eq:depth-sigma} breaks down entirely. A triangulation module must therefore
check its own result before handing it to the filter:
\begin{itemize}
  \item positive depth in every view (cheirality);
  \item sufficient parallax (or baseline relative to depth), or equivalently a well-conditioned Gauss--Newton Hessian;
  \item small reprojection error in every view after refinement;
  \item a plausible depth range.
\end{itemize}
Features that fail are discarded, not fed to the update (AGENTS.md F5, Milestone 10b).

\begin{takeaways}
\begin{itemize}
  \item Rigid correspondences satisfy $\bar x_2\T\sx{t}R\bar x_1 = 0$. $E$ has 5 DOF and unknown scale.
  \item With the rotation from the gyroscope, two correspondences determine the translation direction.
  \item RANSAC iterations grow as $w^{-s}$: prefer minimal solvers. Thresholds in normalized units are $\tau_{px}/f$.
  \item Pure rotation makes $E$ vanish. Detect low parallax and fall back to rotation-only consistency checks.
  \item Triangulate linearly, refine by Gauss--Newton with inverse depth, and reject low-parallax or negative-depth points.
  \item Depth error grows as $Z^2/b$.
\end{itemize}
\end{takeaways}

\begin{exercises}
  \item Verify \cref{thm:epipolar} numerically for random points and a random relative pose, then show that adding an
        outlier breaks it.
  \item Show that $\sx{t}R$ has singular values $(\|t\|, \|t\|, 0)$.
  \item Implement the two-point method with known rotation inside RANSAC on synthetic data with 30\% outliers. Compare the
        number of iterations needed with the five-point method.
  \item Generate a pure-rotation sequence, observe what \code{cv::findEssentialMat} returns, and implement the parallax test.
  \item Show that \eqref{eq:dlt} has rank two per view and that two views with non-parallel rays give a unique solution.
  \item Derive \eqref{eq:depth-sigma}, reproduce \cref{fig:triangulation}, and repeat with three views over the same baseline.
\end{exercises}
